\chapter{O neurônio não é um aparelho só}
% versão completa: chapters/18_eetypes.tex

No primeiro capítulo separamos os neurônios pela substância que liberam. Um engenheiro os separaria de outro jeito: pelo aparelho que cada um deles faz as vezes. E acontece que, sob a mesma aparência, se esconde uma caixa inteira de componentes eletrônicos.

\section*{Uma caixa de componentes}

O balde furado do segundo capítulo é um integrador com vazamento: ele soma os sinais numa janela curta. Com uma janela muito pequena, vira um detector de coincidência. O neurônio que só responde a mudanças e logo se cala é um filtro que deixa passar as variações. O neurônio que trabalha sozinho é um metrônomo, um oscilador; se uma entrada ajusta o seu ritmo, é um oscilador controlado por tensão. O neurônio que emite rajadas de cliques é um gerador de rajadas, que escreve ``traços''. O neurônio que clica no ritmo da onda sonora é um detector de fase.

\pencilfig{18}{\pencilart{18}}{Sob a mesma aparência se esconde uma caixa inteira de componentes: resistor, capacitor, bobina, oscilador.}

Há também aparelhos de que ainda não falamos. O \emph{ressonador} parece um balanço: responde melhor a empurrões de um certo ritmo, como o balanço que só ganha altura se for empurrado na hora certa. Dentro do neurônio, o papel da mola é de uma corrente lenta que resiste às mudanças. O \emph{integrador sem vazamento} guarda um valor por muito tempo, como um capacitor sem fuga --- é assim que funciona a rede que mantém a posição dos olhos quando você olha para o lado. Mas para isso a realimentação precisa ser ajustada quase à perfeição, e esse é um ajuste caro. O \emph{latch} é um neurônio que um empurrão curto coloca em funcionamento prolongado até que o desliguem, como um interruptor com trava; nos neurônios que comandam os músculos, esse modo é ligado pela serotonina.

\section*{Um componente camaleão}

A principal descoberta dessa separação: não são peças diferentes, e sim \emph{modos} diferentes da mesma peça. Um mesmo neurônio pode ser integrador de dia e gerador de rajadas à noite, conforme o que as estações de rádio lhe dizem. Um engenheiro chamaria o cérebro de matriz lógica reprogramável: o circuito é um só, e a função das peças é mudada pelos ajustes.

Que os neurônios sabem trabalhar em todos esses modos está medido. Como o cérebro usa a troca de modos para calcular é, em grande parte, hipótese.

\fullversion{18}
